\section{Eight-arcs over \texorpdfstring{$\F_{25}$}{F25}}\label{sec:q25}

We now take $s=5$, so each empty fixed line carries $N=10$ candidate
pairs.  The profile of an invariant eight-arc is
\[
 (f,e)\in\{(0,4),(2,3),(4,2),(6,1),(8,0)\}.
\]

The four profiles with $f\ne2$ are closed by incidence and moment arguments.
Only $(f,e)=(2,3)$ escapes the structural bounds.  Table~\ref{tab:q25-profiles}
separates these routes; the exceptional row records the exact
kernel-checked minimum and equality classification after normalization.

\begin{table}[ht]
\centering
\small
\begin{tabular}{ccccc l}
\toprule
$f$ & $e$ & $E$ & $M$ & proved $L\ge$ & method \\
\midrule
$0$ & $4$ & $27$ & $12$ & $5$   & secant-index moments \\
$2$ & $3$ & $17$ & $12$ & $32$  & exact minimum and five-orbit classification \\
$4$ & $2$ & $11$ & $10$ & $4$   & fixed-center incidence \\
$6$ & $1$ & $9$  & $6$  & $36$  & uniform criterion \\
$8$ & $0$ & $11$ & $0$  & $110$ & uniform criterion \\
\bottomrule
\end{tabular}
\caption{The five profiles of invariant eight-arcs in $\PG(2,25)$.
Here $L=\Npair(\cC)$.}
\label{tab:q25-profiles}
\end{table}

\subsection*{The four structural profiles}

The following simple center estimate is useful.  A secant orbit joining
two distinct selected conjugate pairs has a fixed center that lies on
neither participating mate line.  At most $f$ fixed-point star lines and
$e-2$ other mate lines through this center are occupied.  Since a fixed
point lies on $s+1$ fixed lines, each such cross-pair secant orbit is
invisible on at least
\begin{equation}\label{eq:cross-invisible}
  s+3-f-e
\end{equation}
empty carriers.

\begin{proposition}[The four nonexceptional profiles]
\label{prop:four-profiles}
Let $\cC$ be an invariant eight-arc in $\PG(2,25)$.
\begin{enumerate}[label=(\alph*),leftmargin=2em]
\item If $f=8$ or $f=6$, then $L\ge110$ or $L\ge36$, respectively.
\item If $f=4$, then $L\ge4$.
\item If $f=0$, then $L\ge5$.
\end{enumerate}
\end{proposition}

\begin{proof}
For $(f,e)=(8,0)$ one has $(E,M)=(11,0)$, and for $(6,1)$ one has
$(E,M)=(9,6)$.  Theorem~\ref{thm:pair-criterion} gives $110$ and $36$.

For $(f,e)=(4,2)$ one has $(E,M,N)=(11,10,10)$.  There are two cross-pair
secant orbits, each invisible on at least two empty carriers by
\eqref{eq:cross-invisible}.  Hence $B\ge4$, and
\eqref{eq:aggregate} gives $L=B+R\ge4$.

It remains to treat $(f,e)=(0,4)$.  Here
\[
  (E,M,N)=(27,12,10).
\]
All twelve nonfixed secant orbits are cross-pair orbits, and
\eqref{eq:cross-invisible} gives $B\ge48$.

We recall the second secant-index moment for an eight-arc:
\begin{equation}\label{eq:second-moment}
  \sum_{x\notin\cC}\binom{r_x}{2}=3\binom84=210,
\end{equation}
where $r_x$ is the number of secants of $\cC$ through $x$.  At fixed
external points, the first secant-index sum is $48$: the four fixed mate
secants contain six fixed points each, contributing $24$ incidences, and
the other $24$ secants each contain their unique fixed orbit center.  Since $r_x\le4$ and
$2\binom n2\le3n$ for $0\le n\le4$, their contribution to
\eqref{eq:second-moment} is at most $72$.

The four occupied fixed carriers are precisely the selected mate lines.
For an external nonfixed point $x$ on one of them, write $r_x=1+u_x$, where
the first secant is the mate line.  Every nonfixed secant meets at most the
two selected mate lines not containing its endpoints, so
$\sum u_x\le48$.  Since $u_x\le3$ and
$\binom{1+u}{2}\le2u$, the occupied-carrier contribution is at most $96$.
Thus the two endpoints of candidates on empty carriers contribute at least
$210-72-96=42$ to \eqref{eq:second-moment}.  Conjugate endpoints contribute
equally.  Moreover, every such endpoint $q$ belongs to a unique candidate
$Q=\{q,\phi(q)\}$ on its unique mate carrier $\ell$, and
$r_q=r_{\phi(q)}=\mu_\ell(Q)$.  Partitioning the endpoint contribution by
these candidates therefore gives
\[
  T:=\sum_{\ell\in\cE}\sum_{Q\in\cQ_\ell}
       \binom{\mu_\ell(Q)}2\ge21.
\]
Because $\mu_\ell(Q)\le4$ and
$\binom n2\le2(n-1)$ for $1\le n\le4$, we have $T\le2R$, whence $R\ge11$.
Finally, \eqref{eq:aggregate} gives
\[
  L=27(10-12)+B+R\ge-54+48+11=5.
\]
\end{proof}

\begin{corollary}[Nonexceptional order-five alternate repair]
\label{cor:q25-nonexceptional-repair}
Let $A$ be an invariant ten-arc in $\PG(2,25)$, let $O\subseteq A$ be a
selected nonfixed conjugate orbit, and put $D=A\setminus O$.  If the fixed
point count of $D$ is $f=0,4,6,8$, then this deletion has at least
$4,3,35,109$ alternate-orbit repairs, respectively.
\end{corollary}

\begin{proof}
The set $D$ is an invariant eight-arc and $O$ is one of its legal
conjugate-pair extensions.  Proposition~\ref{prop:four-profiles} gives the
four total legal-pair bounds $5,4,36,110$.  Removing the one possibility
$O$ gives the asserted alternate counts.
\end{proof}

The remaining profile $(f,e)=(2,3)$ has $(E,M,N)=(17,12,10)$.  The center
estimate gives $B\ge18$, but the aggregate identity would still require a
further collision estimate.  We close this profile by a finite reduction.

\subsection*{The exceptional two-fixed-point profile}

\begin{proposition}[Exact exceptional minimum]\label{prop:f2}
Every invariant eight-arc in $\PG(2,25)$ with exactly two fixed selected
points has at least $32$ fresh legal conjugate-pair extensions.  This bound
is attained.  After the two projective normalizations used below, the
configurations attaining equality form exactly five residual-group orbits,
of sizes
\[
  200,\qquad 400,\qquad 400,\qquad 200,\qquad 400.
\]
Their disjoint union has size $1600$.  Every normalized configuration
outside this union has at least $33$ legal pairs.
\end{proposition}

\begin{proof}[Computer-assisted proof, kernel checked]
Represent $\F_{25}$ as $\F_5[\omega]/(\omega^2-2)$, so
$\phi(\omega)=-\omega$.  We first spell out the reduction to the certified
row set.  Order the two selected fixed points.  Representatives of them are
two independent vectors in $\F_5^3$; extend these to a basis and send that
basis to the standard basis.  The resulting element of
$\mathrm{PGL}(3,5)$ sends the ordered pair to
\[
 A=[0:0:1],\qquad B=[0:1:0].
\]
It commutes with $\phi$.  Any remaining selected point is off the line
$AB$, since the selected set is an arc, and hence has a unique representative
$P=[1:y:z]$.  Write
\[
 y=a+b\omega,\qquad z=c+d\omega
 \quad(a,b,c,d\in\F_5).
\]
The determinants of $(A,P,\phi(P))$ and $(B,P,\phi(P))$ are nonzero scalar
multiples of $b$ and $d$, respectively.  Thus $b,d\ne0$.

Choose any one of the three selected nonfixed conjugate pairs and use one of
its two points as $P$.  The base-field matrix
\[
 T_{y,z}[X:Y:Z]
  =\bigl[X:\,-ab^{-1}X+b^{-1}Y:\,-cd^{-1}X+d^{-1}Z\bigr]
\]
is invertible, fixes $A$ and $B$ projectively, and commutes with $\phi$.
It sends $P$ and $\phi(P)$ to
\[
 [1:\omega:\omega],\qquad[1:-\omega:-\omega].
\]
Changing the ordering of the fixed points, the chosen selected orbit, or its
chosen point changes the normalizing projectivity but not the legal-pair
count.  Thus existence, rather than uniqueness, of this normalization is all
that is used.

Here is the precise correspondence with the finite rows.  The certificate
uses the canonical representatives
\[
 [1:y:z],\qquad [0:1:z],\qquad [0:0:1]
\]
for projective points and a bijective enumeration of the $310$ nonfixed
Frobenius pairs.  The standard pair above has orbit code $5$.  Compatibility
with both $A$ and $B$ forces every selected nonfixed pair to have code at
least $5$; injectivity makes the other two codes strictly larger than $5$.
Ordering those codes gives a unique pair
\[
  6\le b'<c'\le309.
\]
Conversely, the row indexed by $(b',c')$ is the union of $\{A,B\}$, the
standard pair, and the two pairs with codes $b',c'$.  Its determinant
predicate says exactly that every triple in this eight-point set is
noncollinear.  Hence every normalized semantic arc occurs among the
\[
 \binom{304}{2}=46{,}056
\]
rows that the dispatcher checks (rows failing the determinant predicate are
discarded).

The same coordinate dictionary identifies legal extensions, not merely old
configurations.  The orbit-code enumeration is a bijection onto all nonfixed
Frobenius pairs.  In a certified row, its freshness test says that the new
pair is disjoint from the eight old points, and its determinant tests say
that adjoining the pair leaves every triple noncollinear.  These are exactly
the two conditions defining a fresh semantic legal pair.  Consequently the
cardinality computed for a row is $L=\Npair(\cC)$.  Every base-field
projectivity used above bijects the candidate pairs and preserves freshness
and determinants, so this equality and the legal-pair count transport in
both directions.

Here $(E,M,N)=(17,12,10)$, so \eqref{eq:aggregate} reads
$L=B+R-34$ and the target is the coupled inequality $B+R\ge66$.
The finite layer below is load bearing for this coupling; separate bounds
on invisible centers and collision moments do not prove it.

Among the determinant-valid rows, the residual cover has $1{,}189$
class representatives.  Kernel-checked composition certificates
give a lower bound of $32$ on every representative.  Five representatives
attain $32$.  The residual group has order $400$; orbit--stabilizer
certificates give the five displayed orbit sizes, prove their pairwise
disjointness, and show that their union has size $1600$.

Strict certificates place every nonminimizing residual class above $32$,
and an exhaustive dispatcher places every determinant-valid normalized
equality row in the five-orbit union.  The row correspondence and
projective invariance just proved carry the lower bound and equality
exhaustion back to every semantic arc, regardless of the normalization
choices.  Conversely, the same invariance transports the five checked
equalities across their residual orbits.
Thus $32$ is the exact semantic minimum, with the stated complete equality
classification up to normalization.  Section~\ref{sec:formalization}
records the trust boundary and the generated-certificate layers.
\end{proof}

\begin{theorem}[Uniform order-five multiplicity]\label{thm:q25}
Every Frobenius-invariant eight-arc in $\PG(2,25)$ has at least four distinct
fresh nonfixed conjugate-pair extensions.
\end{theorem}

\begin{proof}
The fixed-point count is even and at most eight.  Proposition
\ref{prop:four-profiles} gives respective lower bounds $5,4,36,110$ for
$f=0,4,6,8$, and Proposition~\ref{prop:f2} gives $32$ for $f=2$.
\end{proof}

\begin{corollary}[Uniform order-five alternate repair]
\label{cor:q25-uniform-repair}
Let $A$ be an invariant ten-arc in $\PG(2,25)$ and let $O\subseteq A$ be
any selected nonfixed conjugate orbit.  After deleting $O$, at least one
different legal conjugate orbit repairs the arc; in fact at least three do.
\end{corollary}

\begin{proof}
The remainder $D=A\setminus O$ is an invariant eight-arc, and $O$ is one
of its legal conjugate-pair extensions.  Theorem~\ref{thm:q25} supplies at
least four legal pairs for $D$, so at least three are different from $O$.
\end{proof}

The complete-arc classifications in
\cite{MarcuginiMilaniPambianco2007,CoolsaetSticker2009} show that the
smallest complete arcs in $\PG(2,25)$ have size $12$.  Thus every eight-arc
has an ordinary one-point extension.  That fact does not imply
Theorem~\ref{thm:q25}: ordinary extendability neither makes the legal-point
set Frobenius-stable nor ensures that a legal point and its conjugate are
jointly legal.

\begin{remark}[Scope of the equality classification]\label{rem:census}
Proposition~\ref{prop:f2} classifies equality up to normalization: it says
that some base-field collineation sends an equality case into the certified
$1600$-element union.  It neither asserts uniqueness of that collineation
nor enumerates the semantic arcs attaining $32$.  It makes no equality
claim for the other fixed-point profiles, other arc sizes, or other fields.
\end{remark}
